\chapter{Termo de Abertura do Projeto}

\section{Dados do projeto}

\begin{description}
    \item [Nome do Projeto:] Sistema de Projeto Multimodal para Alunos e Servidores da Universidade de Brasília --- SPMAS-UnB
    \item [Data de abertura:] 02/07/2022
    \item [Código:] E1
    \item [Patrocinador:] Universidade de Brasília
    % \item [Responsável:] ???
    \item [Gerente:] João Luís Takur Baraky Dias \\ 18/0033646 / jlbaraky@aluno.unb.br / +55~(63)~99203~0392
\end{description}

\section{Objetivos}

% O propósito do projeto é desenvolver um sistema de transporte multimodal para alunos e servidores da Universidade de Brasília, com o intuito de facilitar e agilizar a locomoção nos 4 campus da universidade, sendo que este sistema deve ser composto por transportes intra-campi e inter-campi. O projeto deve incluir:
O propósito do projeto é desenvolver um sistema de transporte multimodal para alunos e servidores da Universidade de Brasília, com o intuito de facilitar e agilizar a locomoção nos 4 campus da universidade, sendo que este sistema deve ser composto por transportes intra-campi e inter-campi. O projeto deve incluir:

\begin{itemize}
    \item Estrutura
    \item Forma de controle eletroeletrônico
    \item Interface do usuário
    \item Sistema de abastecimento
\end{itemize}

% Os objetivos a serem efetivados podem ser decompostos em três partes principais (que são os pontos de controle):
Os objetivos a serem efetivados podem ser decompostos em três partes principais (que são os pontos de controle):

\begin{enumerate}
    \item Apresentar na documentação: o resumo, introdução, termo de abertura, equipe de trabalho, análise de mercado e cronograma. Data final: 07/03/2022
    \item Apresentar na documentação o Projeto Conceitual do Produto contendo: características gerais, estrutura, descrição de hardware, análise de consumo energético e descrição de software. Além de melhorar, caso necessário, o que não ficou claro no primeiro ponto de controle. Data final: 04/04/2022
    \item Por fim, a entrega final do relatório deve ser constituída do relatório pronto e de um vídeo de no máximo 15 minutos de duração apresentando o projeto. Data final: 25/04/2022
\end{enumerate}
